\chapter{Writing a backend}
\label{ch:extending}

Rosetta ships nineteen backends, but the generator is \textbf{open for
extension}: you can teach it a new target without editing
\code{rosetta::generate} or any core file. This chapter adds a Lua backend as a
worked example --- Lua because it is small enough to show whole, and because the
real one exists to compare against.

\section{The mental model: two layers}

A backend is two independent pieces that run at two different times, and keeping
them straight is most of the work.

\begin{center}
\small
\begin{tabular}{@{}lL{4.4cm}L{3.9cm}c@{}}
\toprule
\textbf{Layer} & \textbf{File you write} & \textbf{Runs when} &
  \textbf{Reflection?} \\
\midrule
\textbf{Runtime visitor} & \code{rosetta/visitors/lua.h} &
  when the \emph{generated} binding is compiled & \yes \\
\addlinespace
\textbf{Generation backend} & a \code{rosetta::Backend} subclass &
  when the driver runs & \no \\
\bottomrule
\end{tabular}
\end{center}

The split exists because \code{rosetta::generate<Ts\ldots>} \textbf{erases the
type pack up front} into plain strings:

\begin{lstlisting}[style=cpp]
struct GenClass { std::string name; std::string header; std::string doc; };
\end{lstlisting}

So a \code{Backend} never sees a C++ type --- only \code{\{name, header, doc\}}
per class. All the real reflection happens later, in your \emph{visitor}, when
the generated \code{auto\_lua.cpp} is compiled.

\begin{lstlisting}[style=plain]
manifest.json -(rosetta_gen)-> gen/ -(build)-> ./generator -(run)-> bindings/lua/
                                                   |                 auto_lua.cpp
                                                   |                 CMakeLists.txt
                                          LuaBackend::emit()   (compiling auto_lua.cpp
                                                                runs the walk)
\end{lstlisting}

\section{Layer 1: the runtime visitor}

This is where binding actually happens. It implements the \textbf{walk visitor
concept}: three required member templates plus one optional, each receiving a
\code{std::meta::info} of the member and its annotation pack.

\begin{lstlisting}[style=cpp]
// include/rosetta/visitors/lua.h
#pragma once
#include <experimental/meta> // C++ P2995
#include <rosetta/walk.h>    // Rosetta
#include <sol/sol.hpp>       // Lua
#include <string>            // C++ std

namespace rosetta {

    template <typename T> struct LuaVisitor {
        sol::usertype<T> &ut;

        // A data member. `Fld` is the reflected field; `Anns...` is its
        // annotation pack (rosetta::doc / readonly / range / ...).
        template <std::meta::info Fld, auto... Anns> void field(const char *name) {
            using F = [:std::meta::type_of(Fld):];
            if constexpr (ann::has<readonly>(Anns...)) {
                ut[name] = sol::readonly_property(
                    [](const T &s) -> F { return s.[:Fld:]; });
            } else {
                ut[name] = sol::property(
                    [](const T &s) -> F { return s.[:Fld:]; },
                    [](T &s, F v) { s.[:Fld:] = v; });
            }
        }

        template <std::meta::info Fn, auto... Anns> void method_instance(const char *name) {
            ut[name] = &[:Fn:];
        }

        template <std::meta::info Fn, auto... Anns> void method_static(const char *name) {
            ut[name] = &[:Fn:];
        }

        // OPTIONAL: constructors. Omit the method entirely and the walk skips
        // it -- it is detected with `requires`.
        // template <std::meta::info Ctor, auto... Anns> void constructor() { ... }
    };

    // Entry point the generated code calls, one per class.
    template <typename T> void bind_lua(sol::state &lua, const char *name) {
        sol::usertype<T> ut = lua.new_usertype<T>(name);
        LuaVisitor<T> v{ut};
        walk<T>(v);
    }

} // namespace rosetta
\end{lstlisting}

The annotation helpers you will use inside a visitor:

\begin{itemize}
  \item \code{rosetta::ann::has<readonly>(Anns...)} --- is annotation \code{A}
    present?
  \item \code{rosetta::ann::get\_or<doc>(doc\{""\}, Anns...)} --- fetch
    annotation \code{A}, or a fallback.
\end{itemize}

And the reflection splices: \code{[:std::meta::type\_of(Fld):]} for a field's
type, \code{obj.[:Fld:]} to access it, \code{\&[:Fn:]} for a member-function
pointer.

\begin{note}
If your target cannot marshal some parameter or return type, mirror the
\code{napi\_supported} / \code{rest\_supported} pattern --- a \code{consteval}
predicate plus an \code{if constexpr} guard in \code{field} and
\code{method\_*} --- so unbindable members are \emph{skipped} instead of failing
the build. And record them: see Section~\ref{ch:coverage}'s instrumentation
recipe, which is one line per gate.
\end{note}

\section{Layer 2: the generation backend}

This scaffolds the per-target project. It implements:

\begin{lstlisting}[style=cpp]
struct Backend {
    virtual ~Backend() = default;
    virtual void emit(const GenContext &) const = 0;
};
\end{lstlisting}

with everything it needs handed over in \code{GenContext}:

\begin{center}
\small
\begin{tabular}{@{}lL{9.2cm}@{}}
\toprule
\code{out\_dir} & root of the generated tree (\code{bindings/}) \\
\code{lib} & this target's module/library name, from the manifest entry \\
\code{classes} & \code{std::vector<GenClass>} --- \code{\{name, header, doc\}}
  per class \\
\code{user\_include} & absolute directory holding the class headers \\
\code{rosetta\_include} & absolute path to rosetta's \code{include/} \\
\bottomrule
\end{tabular}
\end{center}

A minimal backend writes one \code{.cpp} that includes the visitor and calls
\code{bind\_lua<T>} per class, plus a \code{CMakeLists.txt}:

\begin{lstlisting}[style=cpp]
#include <rosetta/generate.h>
#include <filesystem>
#include <fstream>

struct LuaBackend : rosetta::Backend {
    static void write(const std::filesystem::path &p, const std::string &s) {
        std::filesystem::create_directories(p.parent_path());
        std::ofstream(p) << s;
    }

    void emit(const rosetta::GenContext &c) const override {
        std::string includes, binds;
        for (const auto &k : c.classes) {
            includes += "#include \"" + k.header + "\"\n";
            using rosetta::exposed_of;
            using rosetta::qualified_of;
            binds += "    rosetta::bind_lua<" + qualified_of(k) +
                     ">(lua, \"" + exposed_of(k) + "\");\n";
        }

        const auto dir = c.out_dir / "lua";
        write(dir / "auto_lua.cpp",
              "// Generated by rosetta -- do not edit by hand.\n" + includes +
              "#include <rosetta/visitors/lua.h>\n"
              "#include <sol/sol.hpp>\n\n"
              "extern \"C\" int luaopen_" + c.lib + "(lua_State *L) {\n"
              "    sol::state_view lua(L);\n" + binds +
              "    return 0;\n}\n");

        write(dir / "CMakeLists.txt", /* ... mirroring bindings/python's ... */);
    }
};
\end{lstlisting}

\begin{warn}
Two names, never interchangeable. \code{qualified\_of(k)} gives the
\textbf{C++} spelling --- \code{arch::sinv::Data} --- which is what anything the
generated translation unit compiles must use. \code{exposed\_of(k)} gives the
\textbf{host-language} name: the manifest's \field{expose} override, else the
identifier. Getting these the wrong
way round is how two same-named bound types become ambiguous in the generated
source, and how a rename silently fails to take effect.
\end{warn}

\section{Registering it}

\code{rosetta::generate} dispatches each target through a runtime registry keyed
by the manifest's \field{lang} string:

\begin{lstlisting}[style=cpp]
std::map<std::string, std::shared_ptr<Backend>> &backend_registry();
void register_backend(std::string lang, std::shared_ptr<Backend> backend);
\end{lstlisting}

Register at static-init time with the \code{BackendRegistrar} helper, in one
self-contained plugin translation unit:

\begin{lstlisting}[style=cpp]
// lua_plugin.cpp -- your plugin, owned entirely by you
#include "LuaBackend.h"

static const rosetta::BackendRegistrar
    lua_reg{"lua", std::make_shared<LuaBackend>()};
\end{lstlisting}

Because the registrar runs before \code{main}, the \code{"lua"} backend is
present by the time the driver calls \code{generate} --- \textbf{no edit to
\code{generate} required}. An unregistered target is warned about and skipped.

\begin{note}
Your backend does \emph{not} belong in \code{rosetta::backend}. That namespace
is where rosetta's own implementations live, not a registration point. Put
yours in your own namespace; what you inherit from is \code{rosetta::Backend},
and what you call is the public API below.
\end{note}

\section{Wiring it into a project}

Point the manifest at your plugin with \field{plugins}, and add a target.
\code{rosetta\_gen} compiles each listed source into the generated driver's
\code{add\_executable}, resolving paths relative to the manifest.

\begin{lstlisting}[style=json]
{
  "user_include": "./geom",
  "rosetta_include": "../../include",
  "plugins": ["lua_plugin.cpp"],
  "targets": [
    { "lang": "python", "name": "pygeom" },
    { "lang": "lua",    "name": "mygeom" }
  ],
  "classes": [
    { "name": "Model", "header": "Model.h" },
    { "header": "Point.h" }
  ]
}
\end{lstlisting}

Then the usual flow:

\begin{lstlisting}[style=sh]
../../bin/rosetta_gen manifest.json gen   # gen/ + your plugin in its CMakeLists
cmake -S gen -B gen/build && cmake --build gen/build
./generator bindings                      # emits bindings/python AND bindings/lua
\end{lstlisting}

\section{API reference}

\subsection*{Generation layer (\code{<rosetta/generate.h>})}

\begin{lstlisting}[style=cpp]
struct GenClass { std::string name; std::string header; std::string doc; };

struct GenContext {
    std::filesystem::path out_dir;
    std::string           lib;
    std::vector<GenClass> classes;
    std::string           user_include, rosetta_include;
};

struct Backend {
    virtual void emit(const GenContext &) const = 0;
};

void register_backend(std::string lang, std::shared_ptr<Backend>);

struct BackendRegistrar {
    BackendRegistrar(std::string lang, std::shared_ptr<Backend>);
};

std::map<std::string, std::shared_ptr<Backend>> &backend_registry();
\end{lstlisting}

\subsection*{IR accessors --- how a backend spells a bound type}

\begin{itemize}
  \item \code{qualified\_of(const GenClass\&)} / \code{(const GenEnum\&)} --- the
    \textbf{C++} spelling, for anything the generated TU compiles.
  \item \code{exposed\_of(const GenClass\&)} / \code{(const GenEnum\&)} --- the
    \textbf{host-language} name.
  \item \code{exposed\_object\_of(const GenType\&, const GenContext\&)} --- the
    exposed name of the type an IR entry \emph{references}.
  \item \code{names\_type(const GenType\&, const K\&)} --- does this IR type name
    that bound class or enum?
\end{itemize}

\begin{note}
These six lived in \code{rosetta::gen\_detail} before 2026-08. That spelling
still compiles --- the old names are \code{using}-declarations onto these --- but
it was a ``detail'' namespace documented as an extension point, so they now sit
in \code{rosetta::} with the rest of the backend API.
\end{note}

\subsection*{Runtime layer (\code{<rosetta/walk.h>}) --- the visitor concept}

\begin{itemize}
  \item \code{v.template field<Fld, Anns...>(const char *name);}
  \item \code{v.template method\_instance<Fn, Anns...>(const char *name);}
  \item \code{v.template method\_static<Fn, Anns...>(const char *name);}
  \item \code{v.template constructor<Ctor, Anns...>();} ---
    \textbf{optional}, called only if defined.
  \item \code{rosetta::ann::has<A>(Anns...)} /
    \code{rosetta::ann::get\_or<A>(fallback, Anns...)}
\end{itemize}

\section{Where to copy from}

The built-in backends are a declaration header in
\code{include/rosetta/backends/} with its implementation alongside in
\code{include/rosetta/backends/inline/}. They live in \code{rosetta::backend}
--- \code{backend::Python}, \code{backend::Node} --- so the directory, the file
and the namespace each say ``backend'' once and the class says only which one.
They are the reference for the CMake flags and the clang-p2996 link line.

The reflection-free header-only runtimes those backends emit \code{\#include}
lines for --- the N-API marshaller, the C\#/Java C-ABI shims, the Qt and ImGui
widget kits --- live under \code{include/rosetta/runtime/}. They are compiled by
the \emph{generated} project, never by the generator, which is why they are a
directory of their own.

The existing visitors in \code{include/rosetta/visitors/} are the templates to
copy for layer~1.
